\section{The spot and the front expiry}\label{sec:front}

The spot is not a state variable. By Proposition~\ref{prop:spot} it is the read-off
$S_t=K^{\!\top}q^j_t$, the same scalar at every expiry. What has to be specified is
the front density $q^{1}$, the initial condition of the recursion of
Proposition~\ref{prop:codebook}; everything at $j\ge2$ follows from it and from the
volatility structure.

\subsection{The front chart}\label{sec:frontslice}

The front density must be strictly positive and sum to one. Both are automatic in the
chart
\begin{equation}\label{eq:softmax}
  q^{1,i}=\frac{e^{z^i}}{\sum_ke^{z^k}},\qquad
  \dd z=m\,\dd t+\sum_{f=1}^{d}\nu^f\,\dd W^f ,
\end{equation}
with $z\in\R^{N+2}$ unconstrained. The chart is invariant under $z\mapsto z+c\one$, so a
shift of $\nu$ along $\one$ leaves $q^1$ unchanged and only the $q^1$-centred part of
$\nu$ is identified; write $\mathrm{dev}^f:=\nu^f-(q^1)^{\!\top}\nu^f\,\one$ for it. The input
$\nu$ is free, and plays the role $b$ plays at $j\ge2$. It\^o's formula for a function
of a vector process shows that each $q^{1,i}$ is a martingale exactly when
\begin{equation}\label{eq:frontdrift}
  m^i=-\tfrac12\sum_f\bigl(\mathrm{dev}^{i,f}\bigr)^2
\end{equation}
up to a multiple of $\one$. Furthermore we obtain
\begin{equation}\label{eq:psi1}
  \psi^{1,\cdot,f}=q^1\circ\mathrm{dev}^f,
  \qquad
  \eta^{1,\cdot,f}=\mathrm{dev}^f,
  \qquad
  \one^{\!\top}\psi^{1,\cdot,f}=0 ,
\end{equation}
with $\circ$ the entrywise product.
The same computation appears in \cite{FilipovicHughstonMacrina2012} and as the
stochastic replicator equation \cite{MertikopoulosMoustakas2010}.

\begin{remark}[Resolution at the front expiry]\label{rem:frontres}
Among grid measures with mean $S$ the least dispersed is the pair of nodes straddling
it, so
\begin{equation}\label{eq:varfloor}
  \Var_{q^1}(K)\;\ge\;(S-K^i)(K^{i+1}-S),
  \qquad K^i\le S<K^{i+1},
\end{equation}
which is at most $\tfrac14h^{2}$. The front implied volatility therefore has the floor
\[
  \frac{\sqrt{(S-K^i)(K^{i+1}-S)}}{S\sqrt{T_1-t}}\;\le\;\frac{h}{2S\sqrt{T_1-t}},
\]
largest at the midpoint of a cell and zero when the spot sits on a node, which is
exactly what Proposition~\ref{prop:collapse} forces at $T_1$.
\end{remark}

\subsection{Terminal settlement: pinning the spot to a strike}\label{sec:nonpinning}

The chart constrains $q^1$ at no time $t<T_1$ beyond positivity and mass, so it imposes
no terminal condition. Wissel's front parametrisation fixes the terminal state whatever
the coefficients \cite{Wissel2007}; here the terminal state is a modelling choice. We take
the front volatility proportional to the distance from the spot, in one factor,
\begin{equation}\label{eq:nufront}
  \eta^{1,i}_{t}=\frac{c_t\,(K^i-S_t)}{h},\qquad c_t=a_1\tau^{-p}\ \text{near }T_1,
  \quad\tau=T_1-t,\ a_1>0 .
\end{equation}
Only the behaviour near $T_1$ matters: on $[0,T_1-\varepsilon]$ the front
volatility may be any bounded, locally Lipschitz choice, and \eqref{eq:nufront} need only
hold on $[T_1-\varepsilon,T_1)$. It is already centred, since $\sum_iq^{1,i}(K^i-S_t)=0$. Its coefficients are polynomial in
$q^1$ and bounded on $[0,T]$ for every $T<T_1$, so there is a unique strong solution
on $[0,T_1)$,
\begin{equation}\label{eq:frontexplicit}
  q^{1,i}_{t}=q^{1,i}_{0}\exp\Bigl(\int_0^t\eta^{1,i}\,\dd W-\tfrac12\int_0^t(\eta^{1,i})^2\dd s\Bigr),
  \qquad
  \lVert\eta^{1}_t\rVert_\infty\le c_t\,w/h,\quad w:=K^{N+1}-K^0 .
\end{equation}

\begin{proposition}[The front collapses onto a node, and settles exactly]\label{prop:collapse}
If $p\ge\tfrac12$ then $q^{1}_{T_1}:=\lim_{t\uparrow T_1}q^1_t=e_k$ for a random index $k$,
almost surely, so $S_{T_1}=K^k$ and
\begin{equation}\label{eq:settle}
  \sum_i q^{1,i}_{T_1}(K^i-K^n)^{+}=(S_{T_1}-K^n)^{+}\qquad\text{for every }n .
\end{equation}
If $p<1$ then Assumption~\ref{ass:F} holds on $[0,T_1]$. The window $p\in[\tfrac12,1)$ gives both.
\end{proposition}

\begin{proof}
Each $q^{1,i}$ is a bounded martingale, so $q^1_t\to q^1_{T_1}$ almost surely, and the
variance $V_t:=\sum_iq^{1,i}_{t}(K^i-S_t)^2$ converges to that of $q^1_{T_1}$. By
\eqref{eq:psi1} and \eqref{eq:nufront} the volatility of $S=K^{\!\top}q^1$ is
$\sum_iK^iq^{1,i}\eta^{1,i}=c_tV_t/h$ on $[T_1-\varepsilon,T_1)$. Since $S$ is a martingale with values in
$[K^0,K^{N+1}]$,
\[
  \E\int_{T_1-\varepsilon}^{T_1}\frac{c_t^2V_t^2}{h^2}\,\dd t
  =\E\bigl[(S_{T_1}-S_{T_1-\varepsilon})^2\bigr]\le w^2 ,
\]
so the integral is finite almost surely. On the event $\{V_{T_1}>0\}$ we would have
$V_t\ge V_{T_1}/2$ near $T_1$, and the integral would diverge, because
$\int^{T_1}c_t^2\,\dd t=a_1^2\int^{T_1}\tau^{-2p}\,\dd t=\infty$ for $2p\ge1$. Hence
$V_{T_1}=0$ almost surely: $q^1_{T_1}$ is a probability vector on distinct points with
zero variance, that is $e_k$ with $K^k=S_{T_1}$, and \eqref{eq:settle} reads
$(K^k-K^n)^{+}=(S_{T_1}-K^n)^{+}$. Finally $\eta^1$ is bounded on $[0,T_1-\varepsilon]$ and
$\int_{T_1-\varepsilon}^{T_1}\lVert\eta^1_t\rVert_\infty\dd t\le
(w/h)\int_{T_1-\varepsilon}^{T_1}c_t\,\dd t<\infty$ when $p<1$, by
\eqref{eq:frontexplicit}.
\end{proof}

The terminal spot is therefore a grid point, as in \cite{Wissel2007}; the difference is
that here it is chosen, by \eqref{eq:nufront}, and there it is forced by the
parametrisation.

\begin{remark}[What the lattice could carry, and does not]\label{rem:bary}
At expiry each butterfly price $q^{j,i}$ is worth its payoff: for $S_{T_j}\in[K^k,K^{k+1}]$ the
second difference defining $q^{j,i}$ telescopes to $\varpi(S_{T_j})$, the barycentric
weights of the spot on the two nodes straddling it. Since $u\mapsto(u-K^n)^{+}$ is affine
between nodes,
\begin{equation}\label{eq:barypayoff}
  \sum_i \varpi^i(s)\,(K^i-K^n)^{+}\;=\;(s-K^n)^{+}
  \qquad\text{for every }s\in[K^0,K^{N+1}]\text{ and every model strike }K^n ,
\end{equation}
so $\varpi(s)$ is an admissible terminal state that settles every gridded contract exactly,
also for $s$ strictly inside a cell. The chart does not force the pinning; the choice
\eqref{eq:nufront} does, because its terminal variance must vanish. A front volatility
that drives $q^1$ to $\varpi(s)$ with $s$ inside a cell is left as future work
(Section~\ref{sec:open}).
\end{remark}
